\documentclass[10pt,a4paper]{amsart}
\usepackage{amsmath,amssymb,amsthm,graphicx}
\usepackage{verbatim}
\newenvironment{reference}{\comment}{\endcomment}
%\newenvironment{reference}{}{}
\newenvironment{slogan}{\comment}{\endcomment}
\newenvironment{history}{\comment}{\endcomment}

% For commutative diagrams we use Xy-pic
\usepackage[all]{xy}

% We use 2cell for 2-commutative diagrams.
\xyoption{2cell}
% Theorem environments.
%
\theoremstyle{plain}
\newtheorem{theorem}[subsection]{Theorem}
\newtheorem{proposition}[subsection]{Proposition}
\newtheorem{lemma}[subsection]{Lemma}

\theoremstyle{definition}
\newtheorem{definition}[subsection]{Definition}
\newtheorem{example}[subsection]{Example}
\newtheorem{exercise}[subsection]{Exercise}
\newtheorem{situation}[subsection]{Situation}

\theoremstyle{remark}
\newtheorem{remark}[subsection]{Remark}
\newtheorem{remarks}[subsection]{Remarks}

\numberwithin{equation}{subsection}

% Macros
%
\def\lim{\mathop{\mathrm{lim}}\nolimits}
\def\colim{\mathop{\mathrm{colim}}\nolimits}
\def\Spec{\mathop{\mathrm{Spec}}}
\def\Hom{\mathop{\mathrm{Hom}}\nolimits}
\def\Ext{\mathop{\mathrm{Ext}}\nolimits}
\def\SheafHom{\mathop{\mathcal{H}\!\mathit{om}}\nolimits}
\def\SheafExt{\mathop{\mathcal{E}\!\mathit{xt}}\nolimits}
\def\Sch{\mathit{Sch}}
\def\Mor{\mathop{\mathrm{Mor}}\nolimits}
\def\Ob{\mathop{\mathrm{Ob}}\nolimits}
\def\Sh{\mathop{\mathit{Sh}}\nolimits}
\def\NL{\mathop{N\!L}\nolimits}
\def\CH{\mathop{\mathrm{CH}}\nolimits}
\def\proetale{{pro\text{-}\acute{e}tale}}
\def\etale{{\acute{e}tale}}
\def\QCoh{\mathit{QCoh}}
\def\Ker{\mathop{\mathrm{Ker}}}
\def\Im{\mathop{\mathrm{Im}}}
\def\Coker{\mathop{\mathrm{Coker}}}
\def\Coim{\mathop{\mathrm{Coim}}}

% Boxtimes
%
\DeclareMathSymbol{\boxtimes}{\mathbin}{AMSa}{"02}

%
% Macros for moduli stacks/spaces
%
\def\QCohstack{\mathcal{QC}\!\mathit{oh}}
\def\Cohstack{\mathcal{C}\!\mathit{oh}}
\def\Spacesstack{\mathcal{S}\!\mathit{paces}}
\def\Quotfunctor{\mathrm{Quot}}
\def\Hilbfunctor{\mathrm{Hilb}}
\def\Curvesstack{\mathcal{C}\!\mathit{urves}}
\def\Polarizedstack{\mathcal{P}\!\mathit{olarized}}
\def\Complexesstack{\mathcal{C}\!\mathit{omplexes}}
% \Pic is the operator that assigns to X its picard group, usage \Pic(X)
% \Picardstack_{X/B} denotes the Picard stack of X over B
% \Picardfunctor_{X/B} denotes the Picard functor of X over B
\def\Pic{\mathop{\mathrm{Pic}}\nolimits}
\def\Picardstack{\mathcal{P}\!\mathit{ic}}
\def\Picardfunctor{\mathrm{Pic}}
\def\Deformationcategory{\mathcal{D}\!\mathit{ef}}
\usepackage{cleveref}
\crefname{originalreference}{original reference}{original references}
\Crefname{originalreference}{Original reference}{Original references}
\makeatletter
\newlabel{situation-setup-base-change}{{0FVG}{source}{}{sourcepointer-1}{}}
\newlabel{situation-setup-base-change@cref}{{[originalreference][0][]0FVG}{[1][0][]source}{}{}{}}
\makeatother
\crefname{originallemma}{original lemma}{original lemmas}
\Crefname{originallemma}{Original lemma}{Original lemmas}
\makeatletter
\newlabel{lemma-pullback-coherent}{{02RE}{source}{}{sourcepointer-2}{}}
\newlabel{lemma-pullback-coherent@cref}{{[originallemma][0][]02RE}{[1][0][]source}{}{}{}}
\makeatother
\crefname{originaldefinition}{original definition}{original definitions}
\Crefname{originaldefinition}{Original definition}{Original definitions}
\makeatletter
\newlabel{definition-flat-pullback}{{02RB}{source}{}{sourcepointer-3}{}}
\newlabel{definition-flat-pullback@cref}{{[originaldefinition][0][]02RB}{[1][0][]source}{}{}{}}
\makeatother
\crefname{originallemma}{original lemma}{original lemmas}
\Crefname{originallemma}{Original lemma}{Original lemmas}
\makeatletter
\newlabel{lemma-pullback-coherent-base-change}{{0FVI}{source}{}{sourcepointer-4}{}}
\newlabel{lemma-pullback-coherent-base-change@cref}{{[originallemma][0][]0FVI}{[1][0][]source}{}{}{}}
\makeatother
\title{The flat pullback commutes with the base change pullback}
\author{The Stacks Project Authors}
\date{stacks-project:chow.tex:local-snapshot-20261012; GNU FDL 1.2 or later}
\begin{document}
\maketitle
\noindent\textit{Editorial note.} All delivered mathematical text is the author's own, copied from the pinned Stacks Project chapter file chow.tex; no statement or proof step was rewritten, reordered or removed.\par\medskip
\noindent\textit{Editorial note.} Preamble provenance: the pinned chapter file opens with an input of the shared preamble and carries no preamble text of its own, so the project's single Stacks-wide preamble staged at sources/stacks-chow/preamble.tex is excerpted verbatim as two spans, lines 14-18 (Xy-pic setup) and lines 154-217 (the native theorem declarations and mathematical macros). Nothing was added, removed or redefined; the wrapper class and the amsmath/amssymb/amsthm/graphicx base packages are supplied by the assembler.\par\medskip
\noindent\textit{Editorial note.} The counted statement and proof contain no cross reference, no citation and no included file, so the delivered body is self contained.\par\medskip
\noindent\textit{Editorial note.} Canonical location of the counted proposition: Stacks tag 0FVK (sources/stacks-chow/pinned-original-source.tex line 15302).\par\medskip
\noindent\textit{Allowed standard prerequisites (editorial).}\begin{itemize}
\item The base change pullback on cycles and Chow groups with its degree shift. Same chapter, Lemma labelled lemma-pullback-base-change, invoked in the statement, and the construction of g\textasciicircum{}* at proof lines 15329-15331. Supplies the horizontal maps of the diagrams.
\item The flat pullback on cycles. Same chapter, Definition labelled definition-flat-pullback, cited at proof lines 15333-15335. Gives f\textasciicircum{}*[Z] = [Z x\_X Y]\_\{k+r\}.
\item Base change of the coherent pullback of a cycle. Same chapter, Lemmas labelled lemma-pullback-coherent and lemma-pullback-coherent-base-change, cited at proof lines 15331-15337. Computes the two composites of the cycle level square.
\end{itemize}\medskip
\ifdefined\hypertarget\hypertarget{sourcepointer-1}{}\fi
\noindent\textit{Source reference.} Original reference 0FVG: sources/stacks-chow/pinned-original-source.tex, source lines 15123--15123. Same-chapter pointer into the author environment situation at source line 15123; invoked at proof line None of the counted proof. It is retained only as a navigation pointer; every load-bearing step is inside the delivered span.\par\medskip
\ifdefined\hypertarget\hypertarget{sourcepointer-2}{}\fi
\noindent\textit{Source reference.} Original lemma 02RE: sources/stacks-chow/pinned-original-source.tex, source lines 2331--2331. Same-chapter pointer into the author environment lemma at source line 2331; invoked at proof line 15329 of the counted proof. It is retained only as a navigation pointer; every load-bearing step is inside the delivered span.\par\medskip
\ifdefined\hypertarget\hypertarget{sourcepointer-3}{}\fi
\noindent\textit{Source reference.} Original definition 02RB: sources/stacks-chow/pinned-original-source.tex, source lines 2219--2219. Same-chapter pointer into the author environment definition at source line 2219; invoked at proof line 15332 of the counted proof. It is retained only as a navigation pointer; every load-bearing step is inside the delivered span.\par\medskip
\ifdefined\hypertarget\hypertarget{sourcepointer-4}{}\fi
\noindent\textit{Source reference.} Original lemma 0FVI: sources/stacks-chow/pinned-original-source.tex, source lines 15193--15193. Same-chapter pointer into the author environment lemma at source line 15193; invoked at proof line 15333 of the counted proof. It is retained only as a navigation pointer; every load-bearing step is inside the delivered span.\par\medskip
\begin{lemma}
\label{lemma-pullback-base-change-pullback}
In Situation \ref{situation-setup-base-change} let $Y \to X \to S$ be locally
of finite type and let $Y' \to X' \to S'$ be the base change by $S' \to S$.
Assume $f : Y \to X$ is flat of relative dimension $r$. Then $f' : Y' \to X'$
is flat of relative dimension $r$ and the diagrams
$$
\vcenter{
\xymatrix{
Z_{k + r}(Y) \ar[r]_{g^*} & Z_{k + c + r}(Y') \\
Z_k(X) \ar[r]^{g^*} \ar[u]^{f^*} & Z_{k + c}(X') \ar[u]_{(f')^*}
}
}
\quad\text{and}\quad
\vcenter{
\xymatrix{
\CH_{k + r}(Y) \ar[r]_{g^*} & \CH_{k + c + r}(Y') \\
\CH_k(X) \ar[r]^{g^*} \ar[u]^{f^*} & \CH_{k + c}(X') \ar[u]_{(f')^*}
}
}
$$
of cycle and chow groups commutes.
\end{lemma}

\begin{proof}
It suffices to show the first diagram commutes. To see this, let
$Z \subset X$ be an integral closed subscheme of $\delta$-dimension $k$
and denote $Z' \subset X'$ its base change. By construction
we have $g^*[Z] = [Z']_{k + c}$. By Lemma \ref{lemma-pullback-coherent}
we have $(f')^*g^*[Z] = [Z' \times_{X'} Y']_{k + c + r}$.
Conversely, we have $f^*[Z] = [Z \times_X Y]_{k + r}$ by
Definition \ref{definition-flat-pullback}. By
Lemma \ref{lemma-pullback-coherent-base-change}
we have $g^*f^*[Z] = [(Z \times_X Y)']_{k + r + c}$.
Since $(Z \times_X Y)' = Z' \times_{X'} Y'$ by
associativity of fibre product we conclude.
\end{proof}

\end{document}
